\chapter{Methodology \& Mathematical Formulation}
\label{ch:methodology}

This chapter formulates single-asset trading as an episodic Markov Decision
Process (MDP). It first sets out the state the agent observes, tracing how
that state was arrived at through three successive versions and stating
precisely how far it can be called complete. It then derives, from
first principles, the two priority learning rules of this project --- the
Monte-Carlo policy gradient (REINFORCE) and Deep Q-learning --- and
introduces PPO as a third, stabilised policy-gradient ablation row. Particular
care is taken to show \emph{why} the implemented loss functions follow from
the stated objective --- i.e.\ why the policy-gradient loss contains a
$\log\pi_\theta$ term, and whether each step is exact or a sampling
approximation. The chapter closes with the evaluation protocol used in
Chapter~5.

In plain terms, this chapter answers four questions precisely: what the
agent \emph{sees} (the state), what it may \emph{do} (the actions), what
it is \emph{paid} (the reward), and \emph{how it learns} (the update
rules). Every symbol used is collected in Table~\ref{tab:notation} for
reference.

\begin{table}[h]
\centering
\caption{Notation used throughout Chapters 3--5.}
\label{tab:notation}
\small
\begin{tabular}{ll}
\toprule
\textbf{Symbol} & \textbf{Meaning} \\
\midrule
$s_t \in \mathbb{R}^9$ & state observed at step $t$ (Table~\ref{tab:state}) \\
$a_t \in \{0,1,2\}$    & action: Hold, Buy slice, Sell slice \\
$r_t$                  & reward: one-step wealth change, net of fees \\
$P_t$                  & asset price at day $t$; \; $T$: episode length (days) \\
$C_t,\; h_t$           & cash balance; (fractional) units of the asset held \\
$W_t = C_t + h_t P_t$  & mark-to-market wealth; \; $C_0$: initial cash \\
$c_t,\; v_t$           & cash and exposure as fractions of $C_0$ \\
$\mathrm{dd}_t$        & fall from peak wealth within the episode (rung $S_4$) \\
$c$                    & proportional fee rate (5 basis points $=$ 0.05\%) \\
$\Delta = 0.1\,C_0$    & fixed monetary trade slice \\
$\lambda$              & weight on the drawdown penalty in the risk-aware reward \\
$\gamma$               & discount factor \\
$\tau$                 & a trajectory (one episode); also the clock feature $t/T$ \\
$G(\tau),\; G_t$       & trajectory return; reward-to-go from step $t$ \\
$\pi_\theta(a\mid s)$  & policy network with parameters $\theta$ \\
$Q_\phi(s,a),\;Q_{\phi^-}$ & Q-network with parameters $\phi$; frozen target copy \\
$J(\theta)$            & objective: expected trajectory return \\
$\varepsilon$          & exploration rate ($\varepsilon$-greedy) \\
\bottomrule
\end{tabular}
\end{table}

\section{Overview: trading as a Markov Decision Process}
\label{sec:mdp_overview}

Reinforcement Learning (RL) trains an agent to make sequential decisions by
interacting with an environment. The environment is formalised as a
discrete-time MDP $(\mathcal{S}, \mathcal{A}, P, R, \gamma)$: at step $t$
the agent observes a state $s_t \in \mathcal{S}$, selects an action
$a_t \in \mathcal{A}$, receives a scalar reward $r_t$, and the environment
transitions to $s_{t+1}$ according to dynamics $P(s_{t+1}\mid s_t, a_t)$.
Figure~\ref{fig:rl_loop} shows this interaction loop instantiated for the
trading task.

\begin{figure}[h]
  \centering
  \includegraphics[width=0.82\textwidth]{rl_loop_final}
  \caption{The agent--environment interaction loop for the trading MDP.
  One cycle corresponds to one trading day.}
  \label{fig:rl_loop}
\end{figure}

To keep the formulation concrete, Table~\ref{tab:mdp_mapping} maps the
components of a familiar game-playing task (Flappy Bird, used in our
Phase-0 validation) onto the trading task.

\begin{table}[h]
\centering
\caption{Mapping a popular game MDP to the portfolio trading MDP.}
\label{tab:mdp_mapping}
\small
\begin{tabular}{p{2.4cm}p{5.2cm}p{5.4cm}}
\toprule
\textbf{MDP component} & \textbf{Flappy Bird} & \textbf{Trading (this work)} \\
\midrule
State $s_t$
  & Bird's position and speed, together with the position of the next pipe
  & Market information and the agent's current portfolio position \\
\addlinespace
Action $a_t$
  & Flap or do not flap
  & Hold, Buy, or Sell \\
\addlinespace
Reward $r_t$
  & Survival score based on how long the bird remains in play
  & Change in portfolio wealth after trading fees \\
\addlinespace
Episode $\tau$
  & One complete game session
  & One calendar month of daily trading \\
\addlinespace
Terminal rule
  & The game ends when the bird crashes or hits the ground
  & The episode ends at the final trading day, with any remaining position sold \\
\bottomrule
\end{tabular}
\end{table}

The analogy also fixes the design principle used throughout. In Flappy
Bird the state sees the whole of the world that matters to the decision,
namely the bird's motion and the geometry of the next pipe, and an agent
given that state need never guess. The trading state is held to the same
test. Section~\ref{sec:state} applies it, and shows that the test can be
met in full for one half of the trading world and only approximately for
the other.

\section{Episode definition and data regime}
\label{sec:episode}

\paragraph{Episode.} One episode is one calendar month of daily closing
prices $P_0, P_1, \dots, P_T$ (with $T \approx 18$--$23$ trading days per
month). The agent starts each episode with cash only and is force-liquidated
at the final bar, so no position is carried between episodes.

\paragraph{Independence assumption.} Episodes are treated as mutually
independent, exactly as successive game sessions are independent. This is
a modelling simplification: consecutive months of market data are in fact
correlated (the data are non-IID in time). Handling that temporal
dependence properly belongs to time-series modelling and partially
observable formulations, and is deliberately out of scope; the assumption
and its consequences are revisited in the conclusions.

\paragraph{Why monthly episodes of daily bars?} A single year of data
yields only about 12 monthly episodes but roughly 250 daily transitions.
Eight years (2018--2025) provide 96 months and about 2{,}000 transitions,
split chronologically with no shuffling, so no test information can leak
backwards (Section~\ref{sec:protocol}). Minute-bar episodes
($T{=}390$ per session) are an immediate drop-in upgrade because the
environment interface does not depend on the horizon; this is noted as
future work.

\section{The state: what the agent sees}
\label{sec:state}

\subsection{Two kinds of information}
\label{sec:state_two_halves}

The design principle stated in Section~\ref{sec:mdp_overview} is that the
state should see the whole of the world that matters to the decision, as
the Flappy Bird state does. Before building such a state it is worth
noticing that a trading agent's world has two halves, and that they
behave very differently.

The first half is the agent's own account: how much cash it holds, how
much of the asset it owns, what it paid for that position, and how far
the month has run. The agent controls this half, and none of it is hidden
or estimated. Given the price path, the account's evolution under any
sequence of actions follows exactly from these quantities, by arithmetic
the environment performs in closed form. For this half, complete
observation is not an aspiration but a property that can be checked.

The second half is the market. Here complete observation is impossible in
principle rather than merely inconvenient. Prices are set by the orders,
information and intentions of every other participant, and none of that
appears in a series of closing prices. However many features are computed
from past prices, the future remains only partly predictable from them.
Formally the trading problem is a partially observable MDP
\cite{kaelbling1998}; read from the other direction, it is the classical
observation that prices already reflect the information available about
them \cite{fama1970}.

The Flappy Bird analogy therefore needs one qualification, and it is
better to state it than to be asked for it. Flappy Bird's world is finite
and deterministic, so a handful of numbers genuinely is sufficient. A
market's world is neither. The attainable goal is a state that is
complete on the account side and rich enough on the market side that
adding further features stops changing what the agent does. That second
condition is an empirical claim, and Section~\ref{sec:state_limits} is
explicit about how far the present work can support it.

\subsection{When is a state admissible?}
\label{sec:state_admissible}

It is tempting to design a state by asking which features look useful.
That question is too weak, because it tests only whether the agent can
choose sensibly and says nothing about whether the object being trained is
an MDP at all. Three conditions are required, and they are checked in
order.

\begin{enumerate}
  \item \textbf{Reward measurability.} The reward must be a function of
        the state, the action and the successor state. If it depends on
        anything else, the agent is being scored on information it cannot
        see.
  \item \textbf{Transition closure.} The successor state must be
        computable from the current state, the action and the exogenous
        price, without reaching back into history.
  \item \textbf{Decision sufficiency.} The state must let the agent
        distinguish situations that call for different actions.
\end{enumerate}

Only the third is about usefulness. The first two decide whether the
formulation is legitimate, and they are the reason two apparently similar
quantities are treated very differently below.

\paragraph{Checking the first condition for this work.} The reward of
Section~\ref{sec:reward} is the wealth change $W_{t+1} - W_t$. Since
$W_t = C_t + h_t P_t$, the cash and exposure features of
Table~\ref{tab:state} satisfy $c_t + v_t = W_t / C_0$, and therefore
\begin{equation}
r_t \;=\; C_0\big[(c_{t+1} + v_{t+1}) - (c_t + v_t)\big].
\label{eq:reward_measurable}
\end{equation}
The reward is a function of $(s_t, a_t, s_{t+1})$ with nothing hidden, so
the first condition holds. This is a small calculation but not an idle
one: it is what licenses the derivations later in this chapter.

\paragraph{Two kinds of redundancy.} Equation
(\ref{eq:reward_measurable}) also shows why the episode's running profit,
$W_t/C_0 - 1$, is deliberately \emph{not} a feature. It equals
$c_t + v_t - 1$ exactly, so it is recoverable \emph{pointwise}, from
quantities the agent already observes at the same instant, and a single
linear layer can form it. Supplying it would add width without adding
information.

Drawdown behaves differently. Write
$D_t = \max_{u \le t} W_u - W_t$ for the fall from the highest wealth
reached so far. Two episodes can arrive at identical cash, identical
exposure and identical wealth by different routes: one that rose steadily
to its present value and one that reached a much higher peak and gave
most of it back. The second is in a materially different position, and no
combination of current-step features distinguishes the two, because a
running maximum is a property of the \emph{path}. Drawdown is therefore
not redundant; it is simply absent unless it is supplied.

This gives a rule that decides such cases without appeal to taste:
\emph{exclude features that are pointwise recoverable, and supply
features that are path-dependent and relevant.} It is applied twice in
what follows.

\subsection{The design ladder}
\label{sec:state_rationale}

The state was arrived at by starting from the smallest representation in
which a trading decision can be expressed and asking, at each stage, what
situations the agent still cannot tell apart. Each rung is therefore
justified by a limitation of the one before it rather than by a wish for
more inputs. The progression is one of increasing informational
completeness, not of increasing feature count.

Table~\ref{tab:state_ladder} sets out the ladder. It is a design
argument, and it is stated here before any experiment: Chapter~5 tests
whether each addition actually changes what is learned, and the two
questions are kept apart deliberately, so that a reasoned specification is
not confused with evidence about its consequences.

\begin{table}[h]
\centering
\caption{The state-design ladder. Each rung resolves the limitation named
in the row above it. $\Delta P_t$ is the one-step return, $c_t$ cash and
$v_t$ exposure as fractions of initial capital, $h_t$ the units held,
$\tau_t$ the episode clock and $\mathrm{dd}_t$ the fall from peak wealth.}
\label{tab:state_ladder}
\small
\begin{tabular}{p{1.1cm}p{4.6cm}p{6.2cm}}
\toprule
\textbf{Rung} & \textbf{State} & \textbf{Limitation that motivates the next rung} \\
\midrule
$S_0$
  & $[\Delta P_t]$
  & The agent sees the market but not itself. It cannot tell whether it
    holds a position, or whether a Buy is affordable, so the same market
    observation covers situations needing opposite actions. \\
\addlinespace
$S_1$
  & $[\Delta P_t,\, c_t,\, h_t]$
  & It knows how many units it holds but not what they are worth. Ten
    units is a small position at one price and a large one at another,
    and a feed-forward network cannot form the product $h_t P_t$
    reliably from its inputs. \\
\addlinespace
$S_2$
  & $[\Delta P_t,\, c_t,\, v_t,\, \mathrm{PnL}_t]$
  & It can value its position and tell profit from loss, but not where it
    stands in the month. A Buy on the first day and a Buy the day before
    forced liquidation are different decisions that look identical. \\
\addlinespace
$S_{3a}$
  & $[\Delta P_t,\, c_t,\, v_t,\, \mathrm{PnL}_t,\, \tau_t]$
  & Its own circumstances are now fully described, but its view of the
    market is still a single lagged return. It cannot tell a sustained
    trend from a reversal, or a calm day from a turbulent one. \\
\addlinespace
$S_3$
  & adds trend, volatility, distance from trend, and daily range
    (Table~\ref{tab:state})
  & Complete with respect to the environment as defined, but with no
    memory of the path by which its wealth was reached. \\
\addlinespace
$S_4$
  & adds $\mathrm{dd}_t$
  & Required if the reward refers to risk history
    (Section~\ref{sec:risk_extension}). \\
\bottomrule
\end{tabular}
\end{table}

Two remarks on how to read the ladder. First, the step from $S_1$ to
$S_2$ is a replacement rather than an addition: the raw unit count gives
way to the position's value together with its unrealised return. Moody
and Saffell make the underlying point in arguing that a trading system
must be told its own position, and told it in a form it can use
\cite{moody2001}. Second, the rungs are cumulative from $S_2$ onward, so
any difference in behaviour between neighbouring rungs is attributable to
the information added between them.

\subsection{The state used in this work}
\label{sec:state_def}

Rung $S_3$ is the representation carried into the experiments. It
contains nine features in four blocks: what the \emph{market} is doing,
what the agent's \emph{portfolio} holds, how the current \emph{position}
is performing, and how much \emph{time} remains. Every feature is
scale-free, being a return, a ratio, or a fraction of initial capital, so
the network never sees a raw dollar price. This matters because the price
of the traded instrument roughly doubles over the sample, and a network
fed absolute prices would face inputs from training that never recur at
test time.

\begin{table}[h]
\centering
\caption{The state $s_t \in \mathbb{R}^9$ used in this work (rung $S_3$).
$C_0$ is initial cash, $h_t$ the units held, $W_t = C_t + h_t P_t$ the
mark-to-market wealth. Windows are $k{=}5$ for momentum and volatility,
$m{=}20$ for the trend reference and $n{=}14$ for the average true range.}
\label{tab:state}
\small
\begin{tabular}{cllp{5.1cm}}
\toprule
\# & \textbf{Feature} & \textbf{Definition} & \textbf{What it tells the agent} \\
\midrule
\multicolumn{4}{l}{\emph{Market}} \\
1 & $\Delta P_t$      & $P_t/P_{t-1} - 1$ & the immediate price movement \\
2 & $\mathrm{mom}_t$  & $P_t/P_{t-k} - 1$ & whether that movement is part of a trend \\
3 & $\mathrm{vol}_t$  & std.\ of the last $k$ one-step returns & how variable the market currently is \\
4 & $\mathrm{gap}_t$  & $P_t/\mathrm{MA}_m(P) - 1$ & how far price sits from its recent average \\
5 & $\mathrm{rng}_t$  & $\mathrm{ATR}_n(P)/P_t$ & how wide the daily range is, which closing prices cannot show \\
\addlinespace
\multicolumn{4}{l}{\emph{Portfolio}} \\
6 & $c_t$ & $C_t/C_0$ & how much cash is available to deploy \\
7 & $v_t$ & $h_t P_t/C_0$ & what the open position is worth \\
\addlinespace
\multicolumn{4}{l}{\emph{Position}} \\
8 & $\mathrm{PnL}_t$ & $(P_t - \bar{P}^{\mathrm{entry}}_t)/\bar{P}^{\mathrm{entry}}_t$ & whether the position is in profit or loss \\
\addlinespace
\multicolumn{4}{l}{\emph{Time}} \\
9 & $\tau_t$ & $t/T \in [0,1]$ & how long remains before forced liquidation \\
\bottomrule
\end{tabular}
\end{table}

\paragraph{Why the daily range earns its place.} Features 1 to 4 are all
functions of the closing-price series. Feature 5 is not, and that is
precisely the argument for it. Consider two days that open and close at
the same price: one drifts within a narrow band, the other falls sharply
and recovers by the close. Their closing returns are identical, their
momentum and volatility histories are identical, and any indicator
computed from closes alone must treat them as the same day. They are not
the same day, and the difference is one of risk. Because the average true
range is built from the daily high and low \cite{wilder1978}, it supplies
information that no re-encoding of closing prices can recover. This
distinction, between a feature that adds information and one that merely
re-expresses what is already present, is the test used throughout: a
candidate feature is rejected if it is a function of data the agent
already observes.

\paragraph{A note on how the features are computed.} All market features
are calculated on the continuous price series and only then divided into
monthly episodes. Computing a rolling window inside an episode would
leave the early days of every month with no history to average over, so
the agent would begin each month believing the market had no trend and no
volatility. Momentum does not reset because a new month has started. The
implementation and its verification are described in
Section~\ref{sec:testing}.

\subsection{The risk-aware extension}
\label{sec:risk_extension}

Rung $S_4$ adds a single feature,
\begin{equation}
\mathrm{dd}_t \;=\; \frac{\max_{u \le t} W_u - W_t}{\max_{u \le t} W_u},
\label{eq:drawdown}
\end{equation}
the fraction of its peak wealth the agent has given back within the
current episode. It is kept out of the main state and introduced as an
extension for a reason that is methodological rather than practical. Under
the wealth-change reward, drawdown is not needed: the first admissibility
condition is already satisfied by (\ref{eq:reward_measurable}), and adding
the feature would enlarge the observation to serve a reward that is not
being used. Folding it in would also make its effect impossible to
measure, since there would be no comparison in which it was absent.
Separating it converts an assumption into a question that Chapter~5 can
answer: does explicitly representing risk history change what the agent
learns?

The extension becomes \emph{necessary}, rather than optional, as soon as
the reward refers to risk. A reward of the form
\begin{equation}
r'_t \;=\; \Delta W_t - \lambda\,C_0 \max\!\left(0,\; \mathrm{dd}_{t+1} - \mathrm{dd}_t\right)
\label{eq:risk_reward}
\end{equation}
depends on peak wealth, which $S_3$ does not contain. Pairing
(\ref{eq:risk_reward}) with $S_3$ would violate the first admissibility
condition, and the resulting object would not be an MDP, so the
derivations in this chapter would not apply to it. That combination is
therefore excluded from the experimental design rather than run and
explained afterwards; Section~\ref{sec:protocol} states which
combinations are admissible. The environment enforces the constraint at
construction, so an inadmissible configuration cannot be built by
accident.

\subsection{What this representation does not contain}
\label{sec:state_limits}

By the division of Section~\ref{sec:state_two_halves}, the portfolio,
position and time blocks are complete. Together with the price path, they
determine the account's evolution under any sequence of actions exactly,
and the accounting identity test of Section~\ref{sec:protocol} verifies
this numerically in every episode. The market block is an approximation,
and the honest formulation of the claim is therefore that $s_t$ is
complete \emph{with respect to the trading environment as defined}, and
is a fully observable approximation of a market that is not.

The specific omissions are worth naming. Trend and volatility are
measured over one window length, so a change of regime that shows up as a
shift between short-horizon and long-horizon variability is not directly
represented. Traded volume, which carries information about the conviction
behind a price move, is available in the data but excluded from $S_3$; it
is examined as a further rung rather than assumed to help. Nothing outside
the traded instrument enters the state, so forward-looking measures such
as implied volatility, and any cross-asset or macroeconomic context, are
unavailable.

Other quantities cannot be added at all. Order-book depth and order flow,
the text of news as it arrives, the positions and intentions of other
participants, and the latent state of the market are unobserved by
construction rather than by omission. Th\'eate and Ernst attribute the
overfitting in their own study to an observation space that was too
limited \cite{theate2021}, which is a reminder that this boundary has
practical consequences and not merely formal ones. The standard remedies,
recurrent policies or an explicit belief state, lie outside the scope
agreed for this project and are set out as future work in Chapter~6.

Two scope decisions also bound what the state could usefully contain. The
agent trades a single instrument, whereas a portfolio in the sense of
Markowitz \cite{markowitz1952} requires a vector of weights across assets
and an estimate of how they move together. The agent is also long-only,
so it may decline to participate but cannot express a negative view.

Finally, extending a state is not free. The training period supplies 60
months, which is roughly 1{,}250 daily transitions. Every additional
input increases the number of parameters to be estimated from that fixed
and rather small sample, raising the variance of what is learned and the
risk of fitting one particular history rather than the process behind it
\cite{lopezdeprado2018}; Bellman's observation about the cost of dimension
applies directly \cite{bellman1957}. A state should therefore be extended
only while further features continue to change measured behaviour, which
is why the ladder above is accompanied by the stopping rule declared in
Section~\ref{sec:protocol} rather than left to judgement after the fact.

\section{Action space: a divisible asset}
\label{sec:actions}

The action set is fixed and discrete:
\[
\mathcal{A} = \{\,0:\text{Hold},\quad 1:\text{Buy},\quad 2:\text{Sell}\,\}.
\]

\paragraph{Action representation is a design choice, not a detail.} The
state is not the only thing that has to be represented well. Under a
whole-share model, ``Buy one share'' commits a different amount of money
at every price level, so the economic meaning of the action drifts across
several years of data, and the action becomes infeasible whenever cash
falls below one share price. The same label then denotes a small position
early in the sample and a large one late in it, which is a
representational defect rather than a market fact.

We therefore treat the asset as divisible, as with a fractional-share
broker or a cryptocurrency: a Buy converts a fixed monetary slice
$\Delta = 0.1\,C_0$ of cash into asset units at the current price, and a
Sell liquidates up to the same slice of position value. Ten consecutive
Buys move the portfolio from all-cash to approximately fully invested, and
the meaning of one action is constant throughout.

Because state design and action design are separate representational
axes, changing both at once would make any difference in behaviour
impossible to attribute. The action model is therefore fixed at
$\Delta = 0.1\,C_0$ before the agents are compared, and the whole-share
alternative is retained as a single controlled comparison at the frozen
state (Section~\ref{sec:protocol}) rather than crossed with every rung of
the ladder.

\paragraph{Execution and fees.} With proportional fee rate $c$ (5 basis
points), a Buy spending $x$ dollars acquires $x / (P_t(1+c))$ units, and a
Sell of $u$ units returns $u P_t (1-c)$ dollars.

\paragraph{Feasibility masking.} Buy requires cash above a minimum trade
value; Sell requires a position above the same minimum. Infeasible actions
are masked: during learning and evaluation their probabilities (REINFORCE)
or values (Deep Q) are set to $-\infty$ before the action is chosen, so
illegal trades are never executed. The environment additionally forces
Hold if an illegal action is ever proposed, and counts such proposals as a
standing diagnostic.

\paragraph{Terminal rule.} At the final bar the environment liquidates any
remaining position at the closing price (fee included), so every episode
ends in cash and episode wealth is unambiguous.

\section{Reward function}
\label{sec:reward}

The reward is the one-step change in mark-to-market wealth, net of fees:
\begin{equation}
r_t \;=\; W_{t+1} - W_t, \qquad W_t \;=\; C_t + h_t P_t .
\label{eq:reward}
\end{equation}

Three worked cases make the incentive explicit. Holding cash only:
$r_t = 0$. Holding $h$ units through a price move:
$r_t = h\,(P_{t+1}-P_t)$. Executing a trade of $x$ dollars: the fee
$c\,x$ is paid immediately inside the wealth change, so every unnecessary
round-trip costs $\approx 2cx$. Because rewards telescope,
$\sum_{t=0}^{T-1} r_t = W_T - W_0$: maximising the sum of step rewards is
exactly maximising final wealth. This identity is verified numerically in
every experiment (the ``accounting identity'' test).

\paragraph{A fully numeric example.} Let $C_0 = \$10{,}000$, price
$P_t = \$500$, fee $c = 0.0005$. The agent Buys one slice
($\Delta = \$1{,}000$): it acquires $1000 / (500 \times 1.0005) = 1.999$
units, and its wealth drops immediately by the fee's worth,
$r_t \approx -\$0.50$. Overnight the price rises 1\% to \$505: holding
$1.999$ units, $r_{t+1} = 1.999 \times \$5 \approx +\$10.00$. If it then
Sells the slice at \$505, it receives
$1.999 \times 505 \times 0.9995 \approx \$1{,}009$, paying another
$\approx\$0.50$ fee. Net effect of the round trip: about $+\$9.00$ ---
the \$10 price gain minus \$1 of fees. The same trip on a flat price
loses \$1. The agent is never told any of this; it must discover from
such rewards that trading has a price and only pays when the market moves
favourably.

\paragraph{On the choice of a risk-neutral reward.} Equation
(\ref{eq:reward}) pays for expected wealth and is silent about risk: a
policy that earns \$100 steadily and one that oscillates wildly around
the same mean earn identical total reward. This is a deliberate,
supervision-agreed scoping decision --- it keeps the objective simple
enough to derive and audit end-to-end --- but it has a predictable
consequence, developed in Chapter~5: whenever the market drifts upward on
average, \emph{maximum exposure is optimal}, and a rational learner
should discover buy-and-hold. Risk-aware alternatives (e.g.\
$r_t = \Delta W_t - \lambda D_t$ with a drawdown penalty, or the
differential Sharpe ratio of Moody and Saffell \cite{moody2001}) change
that incentive and are analysed as future work in Chapter~6.

\paragraph{A constraint linking the reward to the state.} Changing the
reward is not a free choice, because the two must remain consistent. Any
quantity the reward depends on must be observable, or the first
admissibility condition of Section~\ref{sec:state_admissible} fails and
the learning rules derived below no longer apply. The drawdown-penalised
reward (\ref{eq:risk_reward}) is exactly such a case, which is why it is
paired only with the extended state $S_4$ and never with $S_3$.

\section{Objective function}
\label{sec:objective}

\subsection{Trajectory return}

A trajectory (episode) is
$\tau = (s_0, a_0, r_0, s_1, \dots, s_T)$, and its discounted return is
\begin{equation}
G(\tau) \;=\; \sum_{t=0}^{T-1} \gamma^{\,t}\, r_t ,
\qquad \gamma \in (0,1].
\label{eq:return}
\end{equation}

\subsection{Where the parameters live}

The agent's decision rule is a \emph{policy network}
$\pi_\theta(a \mid s)$: a neural network with weights $\theta$ that maps a
state to a probability distribution over the three actions. The weights
$\theta$ parameterize the \emph{policy} --- not the reward. The reward
function (\ref{eq:reward}) is fixed by market mechanics; $\theta$
influences the collected rewards only \emph{indirectly}, by changing which
actions are sampled and therefore which trajectories occur. Formally, the
probability of a trajectory under policy $\pi_\theta$ is
\begin{equation}
p_\theta(\tau) \;=\; \rho(s_0)\,
\prod_{t=0}^{T-1} \pi_\theta(a_t \mid s_t)\; P(s_{t+1} \mid s_t, a_t),
\label{eq:traj_prob}
\end{equation}
where $\rho(s_0)$ is the initial-state distribution and $P$ the (unknown)
market dynamics. The objective is the \emph{expected} return under this
distribution:
\begin{equation}
J(\theta) \;=\; \mathbb{E}_{\tau \sim p_\theta}\!\left[\, G(\tau) \,\right],
\qquad
\theta^\ast = \arg\max_\theta J(\theta).
\label{eq:objective}
\end{equation}
Note that $\theta$ enters (\ref{eq:objective}) only through the sampling
distribution $p_\theta$. If it did not --- if nothing in the expression
depended on $\theta$ --- the objective would be a constant and there would
be nothing to optimise.

In plain terms: a single month's profit is luck as much as skill, so we
cannot judge a policy by one episode. The objective $J(\theta)$ is the
\emph{average} profit the policy would earn across all the months it
could face, weighted by how likely each is. Changing the network weights
$\theta$ changes which decisions get made, which changes which futures
occur, which changes the average --- that indirect chain is exactly what
the expectation notation encodes, and what the next section learns to
differentiate.

\paragraph{On the discount factor.} We use $\gamma = 0.99$. With
episodes of only $T \approx 20$ steps, discounting is nearly neutral
($0.99^{20} \approx 0.82$), so the agent is effectively asked to maximise
undiscounted end-of-month wealth; $\gamma < 1$ is retained for numerical
conditioning and consistency with the standard formulation
\cite{sutton2018}.

\section{From the objective to the REINFORCE loss: why the log appears}
\label{sec:pg_derivation}

This section answers, step by step, how the implemented loss
$L(\theta) = -\sum_t \log \pi_\theta(a_t\mid s_t)\, G_t$ follows from
(\ref{eq:objective}).

\subsection{Step 1 --- differentiate under the integral (exact)}

\begin{equation}
\nabla_\theta J(\theta)
= \nabla_\theta \!\int p_\theta(\tau)\, G(\tau)\, d\tau
= \int \nabla_\theta\, p_\theta(\tau)\; G(\tau)\, d\tau .
\end{equation}
This cannot be estimated from samples yet: $\nabla_\theta p_\theta(\tau)$
is not an expectation under $p_\theta$.

\subsection{Step 2 --- the likelihood-ratio (log-derivative) identity (exact)}

For any positive function, calculus gives
$\nabla_\theta \log p_\theta = \nabla_\theta p_\theta / p_\theta$, hence
\begin{equation}
\nabla_\theta\, p_\theta(\tau) \;=\; p_\theta(\tau)\,
\nabla_\theta \log p_\theta(\tau).
\label{eq:logtrick}
\end{equation}
\emph{This identity is the entire reason a logarithm appears in the loss.}
Substituting (\ref{eq:logtrick}):
\begin{equation}
\nabla_\theta J(\theta)
= \int p_\theta(\tau)\, \nabla_\theta \log p_\theta(\tau)\, G(\tau)\, d\tau
= \mathbb{E}_{\tau\sim p_\theta}\!\left[\,
   \nabla_\theta \log p_\theta(\tau)\; G(\tau) \,\right],
\end{equation}
which \emph{is} an expectation under $p_\theta$ and can therefore be
estimated by sampling trajectories from the current policy.

\subsection{Step 3 --- the unknown market dynamics cancel (exact)}

Taking the log of (\ref{eq:traj_prob}),
\[
\log p_\theta(\tau) = \log\rho(s_0)
 + \sum_{t=0}^{T-1}\log \pi_\theta(a_t\mid s_t)
 + \sum_{t=0}^{T-1}\log P(s_{t+1}\mid s_t,a_t),
\]
and only the middle sum depends on $\theta$. The gradient of the dynamics
terms is zero --- which is why the method needs \emph{no model of the
market}:
\begin{equation}
\nabla_\theta J(\theta)
= \mathbb{E}_{\tau\sim p_\theta}\!\left[\;
   \sum_{t=0}^{T-1} \nabla_\theta \log \pi_\theta(a_t \mid s_t)\;
   G(\tau) \;\right].
\label{eq:pg_full}
\end{equation}

\subsection{Step 4 --- causality: reward-to-go (exact in expectation)}

An action at time $t$ cannot influence rewards earned before $t$, and the
corresponding cross-terms have zero expectation. The argument in brief:
for $t' < t$, the term
$\mathbb{E}[\nabla_\theta \log \pi_\theta(a_t \mid s_t)\, r_{t'}]$
factorises by conditioning on the history up to $t$; the inner
expectation
$\mathbb{E}_{a_t \sim \pi_\theta}[\nabla_\theta \log \pi_\theta(a_t \mid s_t)]
= \nabla_\theta \sum_a \pi_\theta(a \mid s_t) = \nabla_\theta 1 = 0$,
because probabilities sum to one whatever $\theta$ is (see
\cite{sutton2018}, ch.~13, for the full treatment). Replacing
$G(\tau)$ by the reward-to-go
$G_t = \sum_{t'=t}^{T-1}\gamma^{\,t'-t} r_{t'}$ therefore leaves the
gradient unbiased while reducing its variance:
\begin{equation}
\nabla_\theta J(\theta)
= \mathbb{E}_{\tau\sim p_\theta}\!\left[\;
   \sum_{t=0}^{T-1} \nabla_\theta \log \pi_\theta(a_t \mid s_t)\; G_t
  \;\right].
\label{eq:pg_togo}
\end{equation}

\subsection{Step 5 --- Monte-Carlo estimation (approximation)}

The expectation is intractable, so it is approximated by the empirical
mean over $N$ sampled trajectories --- this is the only approximation in
the derivation, and it is unbiased:
\begin{equation}
\nabla_\theta J(\theta) \;\approx\; \frac{1}{N} \sum_{n=1}^{N}
  \sum_{t=0}^{T_n-1} \nabla_\theta \log \pi_\theta(a^{(n)}_t \mid s^{(n)}_t)\;
  G^{(n)}_t .
\label{eq:pg_mc}
\end{equation}
The double sum --- over trajectories $n$ and over steps $t$ within each
trajectory --- and the $1/N$ factor are precisely the empirical version of
the expectation in (\ref{eq:objective}).

\subsection{Step 6 --- the surrogate loss implemented in PyTorch}

Automatic differentiation frameworks minimise. Defining
\begin{equation}
L(\theta) \;=\; -\,\frac{1}{N}\sum_{n=1}^{N}\sum_{t=0}^{T_n-1}
 \log \pi_\theta(a^{(n)}_t \mid s^{(n)}_t)\; G^{(n)}_t ,
\label{eq:reinforce_loss}
\end{equation}
where the returns $G^{(n)}_t$ are treated as constants (no gradient flows
through them), gives
$\nabla_\theta L(\theta) = -\nabla_\theta J(\theta)$ evaluated on the
sample. Hence \texttt{loss.backward(); optimizer.step()} performs
stochastic gradient \emph{ascent} on $J$. The minus sign converts
maximisation into the minimisation PyTorch performs; the $\log$ is not a
modelling choice but the direct consequence of the likelihood-ratio
identity in Step~2.

In the implementation the returns within each episode are standardised
(zero mean, unit variance) before entering (\ref{eq:reinforce_loss}); this
is a standard variance-reduction baseline that leaves the gradient
direction unbiased.

\section{Deep Q-learning formulation}
\label{sec:dqn}

Q-learning is the canonical \emph{value-based} alternative: instead of
parameterizing the policy directly, it learns the action-value function.
In plain terms, where REINFORCE learns a \emph{habit} (``in situations
like this, I tend to buy''), Q-learning learns a \emph{price list}
(``in this situation, buying is worth \$3, holding \$1, selling
$-\$2$'') and then simply takes the best-priced action. The price list
obeys a self-consistency condition first written down by Bellman
\cite{bellman1957}:
\begin{equation}
Q^\ast(s,a) = \mathbb{E}\!\left[\, r_t + \gamma \max_{a'} Q^\ast(s_{t+1},a')
 \;\middle|\; s_t = s,\; a_t = a \right],
\label{eq:bellman}
\end{equation}
the expected return of taking action $a$ in state $s$ and acting optimally
thereafter. The greedy policy $a = \arg\max_a Q^\ast(s,a)$ is then optimal.

\paragraph{Why not a table.} Tabular Q-learning stores one value per
state--action pair. Our state is continuous in nine dimensions; a table
over any reasonable discretization explodes combinatorially (and
discretization itself throws away information). Deep Q-learning (DQN)
replaces the table with a neural network $Q_\phi(s,a)$ that generalises
across nearby states.

\paragraph{Learning rule.} Transitions $(s_t,a_t,r_t,s_{t+1})$ are stored
in a replay buffer; minibatches are sampled uniformly to break the strong
temporal correlation of consecutive steps. The temporal-difference target
uses a separate, periodically synchronised \emph{target network}
$Q_{\phi^-}$ for stability:
\begin{equation}
y_t = r_t + \gamma\, (1-d_t) \max_{a' \in \mathcal{A}_{\mathrm{legal}}(s_{t+1})}
      Q_{\phi^-}(s_{t+1}, a'),
\qquad
L(\phi) = \frac{1}{B}\sum_{b=1}^{B} \bigl( Q_\phi(s_b,a_b) - y_b \bigr)^2 ,
\label{eq:dqn_loss}
\end{equation}
where $d_t$ indicates episode termination and the max ranges over
\emph{legal} actions only (feasibility masking applied inside the target).
Exploration uses $\varepsilon$-greedy action selection over legal actions,
with $\varepsilon$ annealed linearly.

\paragraph{Contrast with REINFORCE.} REINFORCE updates once per episode
from complete returns (unbiased, high variance, on-policy); DQN updates
every step from bootstrapped one-step targets (biased during learning, low
variance, off-policy via replay). Comparing the two families on the same
MDP is the core experimental design of Chapter~5.

\section{PPO: the stabilised policy gradient (third ablation row)}
\label{sec:ppo}

REINFORCE and Deep Q are the two priority methods of this project; PPO is
included as a third ablation row because it is the standard modern
policy-gradient method, and because its design is a direct response to a
weakness REINFORCE exhibits transparently in our own Phase-0 results
(CartPole: REINFORCE 292 vs.\ PPO 500).

Two patches take REINFORCE to PPO. First, an \emph{advantage} estimate
replaces the raw return: a learned critic $V_\phi(s)$ provides a baseline,
and $A_t = G_t - V_\phi(s_t)$ rescales each gradient term by how much
better the action was than expected, reducing variance without adding
bias. Second, the policy update is \emph{clipped}. Writing the
probability ratio between the new and old policies as
\begin{equation}
\rho_t(\theta) = \frac{\pi_\theta(a_t \mid s_t)}{\pi_{\theta_{\mathrm{old}}}(a_t \mid s_t)},
\end{equation}
PPO maximises the clipped surrogate
\begin{equation}
L^{\mathrm{CLIP}}(\theta) = \mathbb{E}_t\!\left[
  \min\bigl( \rho_t(\theta)\, A_t,\;
  \mathrm{clip}(\rho_t(\theta),\, 1-\epsilon,\, 1+\epsilon)\, A_t \bigr)
\right],
\label{eq:ppo_clip}
\end{equation}
with $\epsilon = 0.2$. The $\min$ with the clipped term removes any
incentive to move the policy ratio outside $[1-\epsilon, 1+\epsilon]$ in
a single update, preventing the destructively large steps that raw policy
gradients permit. Note that $L^{\mathrm{CLIP}}$ still contains
$\pi_\theta$ through the ratio $\rho_t$ --- the same likelihood-ratio
machinery of Section~\ref{sec:pg_derivation}, evaluated against a frozen
snapshot $\pi_{\theta_{\mathrm{old}}}$ rather than differentiated through
a logarithm directly.

Because the from-scratch requirement of this project is discharged by
REINFORCE and DQN, PPO is used via the reference Stable-Baselines3
implementation, with the same environment, data split, fee, budget and
seeds as the scratch methods.

\section{Algorithms}
\label{sec:algorithms}

\begin{algorithm}[h]
\caption{REINFORCE for the trading MDP (as implemented)}
\label{alg:reinforce}
\begin{algorithmic}[1]
\State initialise policy network $\pi_\theta$
\For{iteration $=1, 2, \dots$ until step budget exhausted}
  \State sample a training month; reset environment ($C_0$ cash, no position)
  \For{$t = 0, \dots, T-1$}
    \State compute legal-action mask; set illegal logits to $-\infty$
    \State sample $a_t \sim \pi_\theta(\cdot \mid s_t)$; execute; store $(\log\pi_\theta(a_t|s_t),\, r_t)$
  \EndFor
  \State compute rewards-to-go $G_t$; standardise within the episode
  \State $L(\theta) \gets -\sum_t \log \pi_\theta(a_t \mid s_t)\, G_t$
  \State one Adam step on $\nabla_\theta L$
\EndFor
\end{algorithmic}
\end{algorithm}

\begin{algorithm}[h]
\caption{Deep Q-learning for the trading MDP (as implemented)}
\label{alg:dqn}
\begin{algorithmic}[1]
\State initialise $Q_\phi$; target $Q_{\phi^-} \gets Q_\phi$; replay buffer $\mathcal{D}$
\For{environment step $=1, 2, \dots$ until step budget exhausted}
  \State with prob.\ $\varepsilon$: random legal action; else $a_t = \arg\max_{a\,\mathrm{legal}} Q_\phi(s_t,a)$
  \State execute $a_t$; store $(s_t, a_t, r_t, s_{t+1}, \mathrm{mask}_{t+1}, d_t)$ in $\mathcal{D}$
  \State sample minibatch; form masked targets $y$ per (\ref{eq:dqn_loss}); Adam step on $L(\phi)$
  \State every $K$ steps: $\phi^- \gets \phi$
\EndFor
\end{algorithmic}
\end{algorithm}

\section{Evaluation protocol}
\label{sec:protocol}

\paragraph{Split.} Chronological and three-way, on calendar boundaries:
60 training months (2018--2022), 12 validation months (2023) and 24
held-out test months (2024--2025). Every design and hyper-parameter
decision is taken on the validation year; the test years are evaluated
once. A two-way split would leave no way to demonstrate that the choice of
learning rate, network width or step budget had not been influenced by the
months used to report results.

\paragraph{Evaluation.} Greedy, deterministic action selection with the
same feasibility mask used in training. Every method sees identical
episodes, identical initial cash ($C_0 = \$10{,}000$) and identical fees
(5 basis points).

\paragraph{Reference strategies.} Four, and the distinction between the
two passive ones matters.
\begin{itemize}
  \item \textbf{B0, do-nothing}: always Hold. Preserves cash exactly, so
        any method below it destroys value.
  \item \textbf{B1a, slice schedule}: buy one slice whenever a Buy is
        legal. This is the best a strategy restricted to the agent's own
        action granularity can do, so comparing against it isolates
        timing skill from any advantage in action representation.
  \item \textbf{B1b, buy-and-hold}: fully invested at the first bar and
        held to forced liquidation. This is the passive benchmark an
        investor would actually use, and the harder of the two to beat in
        a market that drifts upward.
  \item \textbf{B2, random legal}: uniform over legal actions, from a
        single generator threaded through all test months, giving the
        chance floor.
\end{itemize}
Reporting both passive references prevents a familiar error. A schedule
that buys a tenth of capital per step needs ten steps to become fully
invested and so holds well under full exposure for much of a month;
presenting it as buy-and-hold understates the benchmark and flatters
whatever is compared against it.

\paragraph{Metrics.} With monthly returns $R_i = \Delta W_i / C_0$ over
the $M{=}24$ test months:
\begin{itemize}
  \item \textbf{Mean and standard deviation of $\Delta W$}: central
        tendency and month-to-month variability of wealth change, in
        dollars.
  \item \textbf{Trades per month}: a behavioural fingerprint, and with
        fees an economic cost the agent should manage.
  \item \textbf{Mean exposure}: the average fraction of capital at risk,
        without which return comparisons are not comparable.
  \item \textbf{Win rate against B1b}: fraction of months with strictly
        higher $\Delta W$ than buy-and-hold.
  \item \textbf{Monthly Sharpe ratio} \cite{sharpe1994}:
        $\bar{R}/\sigma_R$, return earned per unit of return volatility.
        Interpreted with care, because near-zero exposure shrinks the
        denominator faster than the numerator, so Sharpe is only
        meaningful between strategies of comparable exposure.
  \item \textbf{Intra-episode maximum drawdown}: the largest
        peak-to-trough fall \emph{within} a month, reported as a mean and
        a worst case. This is the loss a trader actually lives through.
  \item \textbf{Drawdown of accumulated monthly profit}: the same measure
        applied to monthly profits accumulated on a constant stake. It
        describes month-to-month consistency and is reported separately,
        because it is a different quantity from the one above and not a
        conservative approximation of it.
\end{itemize}

\paragraph{Admissible configurations.} State and reward cannot be crossed
freely, because one combination is not an MDP. Of the four pairings of
$\{S_3, S_4\}$ with the risk-neutral reward (\ref{eq:reward}) and the
risk-aware reward (\ref{eq:risk_reward}), three are admissible: $S_3$ with
the risk-neutral reward is the baseline; $S_4$ with the risk-neutral
reward is the control that isolates the cost of the extra feature; and
$S_4$ with the risk-aware reward is the treatment. The fourth, $S_3$ with
the risk-aware reward, is excluded by
Section~\ref{sec:state_admissible}. A separate comparison varies the
action model at the frozen state, so that state, reward and action
representation are each varied one at a time.

\paragraph{Stopping rule for the ladder.} Declared before any run: the
ladder of Section~\ref{sec:state_rationale} terminates at the first rung
where the next rung changes mean $\Delta W$, intra-episode maximum
drawdown and trades per month by less than the spread observed across
seeds $\{42, 43, 44\}$. Fixing the criterion in advance makes the stopping
point a result rather than a consequence of the time available.

\paragraph{Standing diagnostics.} (i) the accounting identity
$\sum_t \Delta W_t = W_T - W_0$ in every episode, checked on the wealth
change rather than the reward, since under (\ref{eq:risk_reward}) the two
differ by the penalty; (ii) the illegal-action proposal rate; (iii) all
learned results reported over seeds $\{42, 43, 44\}$.

\paragraph{Chapter summary (plain language).} The trading problem is now
fully specified. Each month the agent starts with \$10,000 in cash and
sees nine numbers a human day-trader could also work out: how the price
just moved, whether that move is part of a trend, how choppy the market
is, how far price sits from its recent average, how wide the day's range
was, how much cash it holds, what its position is worth, whether that
position is up or down, and how many days remain. It may hold, buy a
\$1,000 slice, or sell one; it pays a small fee on every trade; and it is
scored by the change in its total wealth.

Those nine numbers were not chosen by taste. Each one answers a situation
the previous version of the state could not tell apart, and one candidate
was deliberately left out because arithmetic already available to the
network could reconstruct it. A tenth number, how far wealth has fallen
from its high point, is held back for a separate experiment, because it is
only needed once the score itself starts penalising risk.

Two learning rules were then derived from that score: REINFORCE, which
strengthens whatever decisions preceded good months, and where the
logarithm in the formula is forced by the calculus rather than chosen; and
Deep Q-learning, which learns a running price list for the three choices.
PPO, a modern stabilised variant, joins as a third contestant. The next
chapter turns these equations into tested code.
